We investigate how foreground mode removal affects the redshift evolution of the H\,I-weighted genus and its response to Alcock--Paczy\'nski (AP) distortions. Using four Jiutian 1G snapshots at $0.10\le z\le0.59$, we construct H\,I mass-weighted fields and measure their genus at a Gaussian smoothing scale of $20\,h^{-1}\mathrm{Mpc}$ under idealized horizontal and wedge cuts. The horizontal cut raises the fitted genus amplitude by a factor of about 4.6, despite reducing the field variance. In the simulation cosmology, the curves show modest evolution over the sampled redshifts; their endpoint changes do not describe the full intermediate variation. Subtracting this intrinsic evolution leaves redshift-dependent residuals under the tested incorrect distance--redshift relations, showing that the filtered measurement retains a response to the assumed geometry. The resulting parameter regions depend strongly on their numerical construction. Tests identify discontinuities associated with integer mesh changes and sharp Fourier masks, while parameter-space smoothing changes the relative region areas. We therefore interpret these regions as exploratory summaries of the simulated pipeline response, rather than calibrated cosmological precision. The results motivate further tests of H\,I topology with numerically stable measurements and realistic foreground models.
